%% ma: L12-B02-0018
%% lop: 12
%% chuong: 1
%% bai: 2
%% dang: 12.2.9
%% loai: TN4
%% muc: TH
%% dap_an: "A"
%% nguon: "Ảnh giáo viên gửi 05/11/2026 (M11): Kiểm tra giữa HK1 Toán 12 (2026-2027), THPT Hồ Thị Bi (TP.HCM), mã đề 1133, Phần I Câu 11"
%% kien_thuc: "Bài 2 (giá trị lớn nhất, nhỏ nhất của hàm số trên đoạn đọc từ đồ thị)"
%% trang_thai: chinh-thuc
%% ghi_chu: "Đề gốc không kèm đáp án; đã giải độc lập. Hình vẽ lại theo hình gốc bằng đường cong trơn qua các điểm (-1;-2), (0;1), (2;-2), (3;1), (4;3), (5;1) (không phải đồ thị của một đa thức cụ thể)."
\begin{ex}
Cho hàm số $f(x)$ liên tục $\mathbb R$ và có đồ thị như hình vẽ sau:
\begin{center}
\begin{tikzpicture}[scale=.8]
  \begin{scope}
    \clip (-1.6,-3.4) rectangle (6,3.9);
    \draw[thick] plot[smooth] coordinates {(-1.300,-3.200)(-1.279,-3.087)(-1.257,-2.984)(-1.236,-2.889)(-1.214,-2.800)(-1.193,-2.717)(-1.171,-2.638)(-1.150,-2.562)(-1.129,-2.488)(-1.107,-2.414)(-1.086,-2.338)(-1.064,-2.261)(-1.043,-2.179)(-1.021,-2.093)(-1.000,-2.000)(-0.929,-1.679)(-0.857,-1.362)(-0.786,-1.050)(-0.714,-0.749)(-0.643,-0.461)(-0.571,-0.190)(-0.500,0.062)(-0.429,0.292)(-0.357,0.494)(-0.286,0.668)(-0.214,0.808)(-0.143,0.913)(-0.071,0.978)(0.000,1.000)(0.143,0.956)(0.286,0.834)(0.429,0.646)(0.571,0.405)(0.714,0.125)(0.857,-0.181)(1.000,-0.500)(1.143,-0.819)(1.286,-1.125)(1.429,-1.405)(1.571,-1.646)(1.714,-1.834)(1.857,-1.956)(2.000,-2.000)(2.071,-1.973)(2.143,-1.895)(2.214,-1.772)(2.286,-1.609)(2.357,-1.412)(2.429,-1.187)(2.500,-0.938)(2.571,-0.671)(2.643,-0.391)(2.714,-0.105)(2.786,0.183)(2.857,0.466)(2.929,0.741)(3.000,1.000)(3.071,1.245)(3.143,1.478)(3.214,1.699)(3.286,1.907)(3.357,2.100)(3.429,2.277)(3.500,2.438)(3.571,2.580)(3.643,2.704)(3.714,2.808)(3.786,2.890)(3.857,2.950)(3.929,2.987)(4.000,3.000)(4.071,2.987)(4.143,2.950)(4.214,2.890)(4.286,2.808)(4.357,2.704)(4.429,2.580)(4.500,2.438)(4.571,2.277)(4.643,2.100)(4.714,1.907)(4.786,1.699)(4.857,1.478)(4.929,1.245)(5.000,1.000)(5.029,0.900)(5.057,0.801)(5.086,0.703)(5.114,0.603)(5.143,0.504)(5.171,0.403)(5.200,0.300)(5.229,0.195)(5.257,0.088)(5.286,-0.022)(5.314,-0.135)(5.343,-0.252)(5.371,-0.374)(5.400,-0.500)};
  \end{scope}
  \draw[truc] (-1.7,0)--(6.1,0) node[below]{$x$};
  \draw[truc] (0,-3.5)--(0,4) node[left]{$y$};
  \draw[phu] (-1,0)--(-1,-2)--(2,-2)--(2,0) (0,1)--(5,1) (3,1)--(3,0) (5,1)--(5,0) (0,3)--(4,3)--(4,0);
  \path (0,0) node[below right]{\footnotesize$O$} (-1,0) node[above left]{\footnotesize$-1$}
        (2,0) node[above right]{\footnotesize$2$} (3,0) node[below]{\footnotesize$3$}
        (4,0) node[below]{\footnotesize$4$} (5,0) node[below]{\footnotesize$5$}
        (0,3) node[above left]{\footnotesize$3$} (0,1) node[above left]{\footnotesize$1$} (0,-2) node[below left]{\footnotesize$-2$};
\end{tikzpicture}
\end{center}
Gọi $A$, $a$ lần lượt là giá trị lớn nhất và giá trị nhỏ nhất của hàm số đã cho trên đoạn $[-1;5]$. Khi đó, giá trị của biểu thức $Q=A-a$ là
\choice{\True $5$}{$1$}{$6$}{$3$}
\loigiai{Từ đồ thị, trên đoạn $[-1;5]$ hàm số đạt giá trị lớn nhất $A=f(4)=3$.
Giá trị nhỏ nhất $a=f(-1)=f(2)=-2$.
Vậy $Q=A-a=3-(-2)=5$.}
\end{ex}
